\section{Experimental Setup}
\label{sec:setup}

\subsection{Benchmark Cases}
We use the two wind data sets of the Kusiak--Song benchmark~\cite{Kusiak2010,Bansal2017}, listed in Table~\ref{S-tab:S-wind} of the supplementary material. In Wind Data Set~I, the Weibull parameters are the same in all directions ($\zeta=2$, $\psi=13$) and the wind is strongly directional (probabilities 0.2 and 0.6 for the bins $[75^\circ,90^\circ)$ and $[90^\circ,105^\circ)$, 0.01 for most others); Wind Data Set~II keeps $\zeta=2$, but $\psi$ varies from 2.6 to 10~m/s with the direction and the wind is spread over more directions. We use $N_\theta=24$ direction bins of $15^\circ$ and $N_s=21$ speed intervals of 0.5~m/s. Each data set is combined with three circular farms of radius $r=500$, 750 and 1000~m, with $N=2,\ldots,10$, $2,\ldots,12$ and $2,\ldots,15$ turbines, respectively, which gives $2\times(9+11+14)=68$ cases. The turbines have $R=38.5$~m, $\ell_{\min}=4D=308$~m, $C_T=0.8$ and $K=0.075$.

\subsection{Methods Compared}
The main comparison includes PSO-VNS and seven reference methods, all evaluated with the objective~\eqref{eq:penalty}, the bounds $[-r,r]^{2N}$, population 30 where applicable, and the same budget:
\begin{itemize}
\item \emph{PSO-VNS} (proposed): Algorithm~\ref{alg:hybrid} with $\omega=0.5$, i.e., 100 PSO iterations (3,030 calls) followed by 3,000 calls of VNS.
\item \emph{PSO}~\cite{Kennedy1995,Clerc2002}: the coefficients \eqref{eq:constriction} and 200 iterations. The setting $w=0.7$, $c_1=c_2=2$ of the earlier version of this study is discussed in Section~\ref{sec:psosetting}.
\item \emph{SSA-VNS}: Algorithm~\ref{alg:hybrid} with SSA in Phase~1 (100 SSA iterations, 3,030 calls, $\xi_1$ decreasing over these 100 iterations).
\item \emph{SSA}~\cite{Mirjalili2017} (200 iterations), \emph{LX-SSA}~\cite{Solanki2023} ($\phi=0$, $\chi=1$, 100 iterations) and \emph{DE}~\cite{Storn1997} (DE/rand/1/bin, scale factor 0.5, crossover rate 0.9, 200 iterations).
\item \emph{VNS}: Phase~2 of Algorithm~\ref{alg:hybrid}, started from the best member of the initial population.
\item \emph{Multistart SLSQP (MS-SLSQP)}: SciPy's SLSQP~\cite{Kraft1988} with explicit spacing and boundary constraints and forward-difference gradients (1-m step), started from the initial population in the order of the $F_p$ values and restarted until the budget is used up.
\end{itemize}
The ablation adds two variants. \emph{RS-VNS} is a random-multistart control: Phase~1 is replaced by pure random sampling with the same number of calls ($\operatorname{round}(\omega B)=3{,}015$ layouts, including the initial population), and the best sampled layout is refined by the same VNS. \emph{LX-SSA-VNS} uses LX-SSA in Phase~1 ($T_1=\operatorname{round}((\omega B-N_p)/(2N_p))=50$ iterations, 3,030 calls). SSA, LX-SSA, PSO and DE use our original implementation from the earlier version of this study, with the coefficients \eqref{eq:constriction} for PSO; VNS, RS-VNS, MS-SLSQP and the three hybrids were implemented for this study, and the hybrids call the original swarm code in Phase~1.

Every run of the main comparison has a budget of exactly 6,030 objective calls, counted by a wrapper around the objective function (including the initial population and, for MS-SLSQP, the finite-difference gradient calls). Each method is run with the seeds 1--30; every optimizer seeds the random-number generator and draws its initial population first, so runs with the same seed start from the \emph{same} initial population for all methods and are paired.

\subsection{Tests Beyond the Benchmark}
\label{sec:setup-beyond}
\paragraph*{Initialization and budget} In the main comparison, the initial layouts are drawn uniformly in the bounding square and are usually infeasible. As an alternative, each initial layout is drawn inside the site and pushed apart by minimizing a smooth packing penalty with L-BFGS-B until the constraints hold, without objective calls. The nine methods other than RS-VNS are run with this initialization at 6,030 calls (for RS-VNS, feasible sampling would require one packing solve for each of its 3,015 random layouts per run, which is far more expensive than the objective calls themselves), and all ten methods (the eight of the main comparison, RS-VNS and LX-SSA-VNS) are run with random initialization at 6,030, 30,030 and 120,030 calls, on the six largest cases (the largest $N$ of each farm and data set) and the Horns Rev~1 block, with 30 seeds (10 for Horns Rev~1 at 120,030 calls). All iteration counts scale with the budget; e.g., PSO-VNS uses 500 PSO iterations at 30,030 calls.

\paragraph*{Horns Rev~1 16-turbine block} We use the Horns Rev~1 offshore site from PyWake~\cite{PyWake}: Vestas V80 turbines (rotor diameter 80~m, 2~MW) with tabulated power and thrust curves and a 25~m/s cut-out, the measured 12-sector wind climate, the parallelogram through the corner turbines of the block as boundary, the Jensen wake with $K=0.04$ (AEP within 0.3\% of PyWake's), and PyWake's 16-turbine example block with $4D$ (320~m) spacing.

\paragraph*{IEA37 Case Study~1} Case Study~1 of IEA Wind Task~37~\cite{Baker2019,IEA37repo} fixes the wake model and leaves the optimizer free; we use its 16- and 36-turbine scenarios (circular boundaries of radius 1300 and 2000~m, minimum spacing $2D=260$~m, IEA 3.35-MW reference turbine with $D=130$~m, 16 wind directions with a constant speed of 9.8~m/s, simplified Bastankhah Gaussian wake). Our implementation reproduces the official AEP calculator to a relative deviation below $10^{-11}$. All methods are run with the penalty~\eqref{eq:penalty} applied to the AEP loss, at 6,030 and 30,030 calls with 30 seeds, and compared with the submitted layouts~\cite{Baker2019,IEA37repo}.

\subsection{Performance Measures and Statistical Analysis}
For each run we record the final benchmark objective, the wake loss (as a percentage of the wake-free objective), whether the final layout is feasible (spacing and boundary satisfied within $10^{-6}$~m), the convergence curve (best feasible objective at 201 equally spaced checkpoints), and the wall-clock time. Means and standard deviations are computed over the feasible runs. Within each case, the 30 seed-paired runs are compared by two-sided Wilcoxon signed-rank tests of PSO-VNS against each other method, with Holm's correction~\cite{Holm1979} over the comparisons of the case and matched-pairs rank-biserial correlations $r_{\rm rb}$ as effect sizes; a feasible run beats an infeasible one. Across the 68 cases, the methods are ranked within each case by the mean benchmark objective of their feasible runs and compared by a Friedman test with the Iman--Davenport correction and Holm-adjusted post hoc tests on the average ranks~\cite{Demsar2006}. Because mean-rank post hoc tests depend on the set of compared methods~\cite{Benavoli2016}, we also compare PSO-VNS with each method by Wilcoxon signed-rank tests on the per-case mean wake losses.

\subsection{Implementation Verification and Environment}
\label{sec:validation}
Our vectorized objective agrees with the original implementation within a relative difference of $10^{-6}$ and reproduces the objective values of 720 archived final layouts within $8.7\times10^{-11}$; re-running the original SSA, LX-SSA, PSO and DE code with the archived settings reproduces the archived result distributions (22 of 24 Mann--Whitney comparisons with $p\ge0.05$). All new runs used Linux containers (Intel Xeon processor at 2.1~GHz, 4 virtual cores) with Python~3.11.15, NumPy~2.4.6 and SciPy~1.17.1.
% CHECK (manual): CPU from lscpu of the analysis container; confirm for the machines that produced all result files


